\documentclass[10pt, aspectratio=169]{beamer}
% Preámbulo modular para presentaciones Beamer (Modelación Estadística)
% Diseño y paleta institucional del Tecnológico de Monterrey

\documentclass[aspectratio=169,xcolor={dvipsnames,table}]{beamer}

\usepackage[utf8]{inputenc}
\usepackage[spanish,mexico]{babel}
\usepackage{amsmath,amssymb,amsfonts,mathtools}
\usepackage{graphicx}
\usepackage{listings}
\usepackage{booktabs}
\usepackage{tikz}
\usetikzlibrary{matrix,arrows,decorations.pathmorphing,shapes.geometric,positioning}

% Paleta Institucional Tecnológico de Monterrey
\definecolor{TecAzul}{HTML}{0039A6}       % Azul TEC: color primario institucional
\definecolor{TecAzulOscuro}{HTML}{1A2E51} % Azul marino oscuro
\definecolor{TecRojo}{HTML}{EC2661}       % Rojo de acento (paleta miTec)
\definecolor{TecGris}{HTML}{646464}       % Gris neutro para comentarios y subtítulos
\definecolor{TecFondoBloque}{HTML}{F4F6F9}% Fondo suave para bloques

% Configuración de Tema Beamer
\usetheme{Boadilla}
\usecolortheme[named=TecAzulOscuro]{structure}
\setbeamercolor{alerted text}{fg=TecRojo}
\setbeamercolor{palette primary}{bg=TecAzulOscuro,fg=white}
\setbeamercolor{palette secondary}{bg=TecAzul,fg=white}
\setbeamercolor{palette tertiary}{bg=TecAzulOscuro!80!black,fg=white}
\setbeamercolor{title}{fg=TecAzulOscuro,bg=TecFondoBloque}
\setbeamercolor{frametitle}{fg=TecAzulOscuro,bg=TecFondoBloque}

% Bloques personalizados elegantes
\setbeamercolor{block title}{bg=TecAzulOscuro,fg=white}
\setbeamercolor{block body}{bg=TecFondoBloque,fg=black}
\setbeamercolor{block title alerted}{bg=TecRojo,fg=white}
\setbeamercolor{block body alerted}{bg=TecRojo!8,fg=black}
\setbeamercolor{block title example}{bg=TecAzul,fg=white}
\setbeamercolor{block body example}{bg=TecAzul!8,fg=black}

% Redondear bordes y añadir sombra sutil
\setbeamertemplate{blocks}[rounded][shadow=true]
\setbeamertemplate{navigation symbols}{}

% Pie de página limpio con número de diapositiva
\setbeamertemplate{footline}{
  \leavevmode%
  \hbox{%
  \begin{beamercolorbox}[wd=.7\paperwidth,ht=2.25ex,dp=1ex,left]{author in head/foot}%
    \hspace*{2ex}\insertshortauthor~~(\insertshortinstitute) \hfill \insertshorttitle
  \end{beamercolorbox}%
  \begin{beamercolorbox}[wd=.3\paperwidth,ht=2.25ex,dp=1ex,right]{date in head/foot}%
    \insertshortdate{}\hspace*{2em}
    \insertframenumber{} / \inserttotalframenumber\hspace*{2ex}
  \end{beamercolorbox}}%
  \vskip0pt%
}

% Configuración de Listings de Python para visualización pedagógica de código
\lstset{
    language=Python,
    basicstyle=\ttfamily\footnotesize,
    keywordstyle=\color{TecAzulOscuro}\bfseries,
    stringstyle=\color{TecRojo},
    commentstyle=\color{TecGris}\itshape,
    showstringspaces=false,
    breaklines=true,
    frame=lines,
    rulecolor=\color{TecAzulOscuro!30},
    backgroundcolor=\color{TecFondoBloque},
    aboveskip=0.5em,
    belowskip=0.5em,
    literate={á}{{\'a}}1 {é}{{\'e}}1 {í}{{\'i}}1 {ó}{{\'o}}1 {ú}{{\'u}}1
             {Á}{{\'A}}1 {É}{{\'E}}1 {Í}{{\'I}}1 {Ó}{{\'O}}1 {Ú}{{\'U}}1
             {ñ}{{\~n}}1 {Ñ}{{\~N}}1 {¿}{{?`}}1 {¡}{{!`}}1
}

% Metadatos institucionales por defecto
\title[Modelación Estadística]{Modelación Estadística para Ciencia de Datos e Ingeniería}
\author[J. Castillo Colmenares]{Juliho David Castillo Colmenares}
\institute[Tec de Monterrey]{Tecnológico de Monterrey \\ Escuela de Ingeniería y Ciencias}
\date{\today}

% Comandos y macros estadísticas compartidas para las presentaciones Beamer (Metropolis)

% Conjuntos numéricos básicos
\newcommand{\R}{\mathbb{R}}
\newcommand{\N}{\mathbb{N}}
\newcommand{\Q}{\mathbb{Q}}
\newcommand{\Z}{\mathbb{Z}}
\newcommand{\set}[1]{\left\{ #1 \right\}}
\newcommand{\abs}[1]{\left|#1\right|}
\newcommand{\norm}[1]{\left\| #1 \right\|}

% Comandos estadísticos y probabilísticos del curso
\newcommand{\Var}{\operatorname{Var}}
\newcommand{\cov}{\operatorname{Cov}}
\newcommand{\corr}{\rho}
\newcommand{\comb}[2]{\begin{pmatrix} #1 \\ #2 \end{pmatrix}}
\newcommand{\s}{\sigma}
\newcommand{\particion}{\mathfrak{P}}
\newcommand{\card}[1]{\#\left( #1 \right)}
\newcommand{\seq}[3]{\set{#1}_{#2}^{#3}}
\newcommand{\E}{\mathbb{E}}
\newcommand{\Prob}{\mathbb{P}}

% Entornos pedagógicos y cajas en español académico natural
\newenvironment{discusion}{%
  \begin{alertblock}{Pregunta de Discusión}%
}{%
  \end{alertblock}%
}

\newenvironment{reflexion}{%
  \begin{alertblock}{Reflexión para el Aula}%
}{%
  \end{alertblock}%
}

\newenvironment{paradoja}{%
  \begin{alertblock}{Paradoja de Exploración}%
}{%
  \end{alertblock}%
}

\newenvironment{motivacion}{%
  \begin{exampleblock}{Motivación: Ciencia de Datos en la Práctica}%
}{%
  \end{exampleblock}%
}

\newenvironment{retoclase}{%
  \begin{block}{Reto Rápido en Equipo}%
}{%
  \end{block}%
}


\title[Distribución Multinomial]{Sección 03.08: Distribución Multinomial}
\subtitle{Ensayos Politómicos, Covarianzas Negativas y Tablas de Contingencia}
\author[J. Castillo Colmenares]{Juliho Castillo Colmenares}
\institute{Tecnológico de Monterrey}
\date{\vspace{-1.2cm}}

\begin{document}

% Slide 1: Portada
\begin{frame}[plain]
  \titlepage
\end{frame}

% Slide 2: Hoja de Ruta
\begin{frame}{Hoja de Ruta --- Capítulo 03: Variables Aleatorias Discretas}
  \footnotesize
  ¿Dónde nos encontramos en el desarrollo modular de la teoría? \pause
  \begin{itemize}\itemsep=0.08em
    \item \textbf{03.01} Concepto, Notación y Definición de Variable Aleatoria \pause
    \item \textbf{03.02} Funciones de Probabilidad (PMF) y Distribución Acumulada (CDF) \pause
    \item \textbf{03.03} Esperanza Matemática y Varianza Operativa \pause
    \item \textbf{03.04} Distribución Binomial y Ensayos de Bernoulli \pause
    \item \textbf{03.05} Distribuciones Geométrica y Binomial Negativa \pause
    \item \textbf{03.06} Distribución Hipergeométrica \pause
    \item \textbf{03.07} Distribución de Poisson y Procesos de Llegada \pause
    \item \color{TecRojo}\textbf{03.08 Distribución Multinomial \emph{(Hoy)}}\color{black}
  \end{itemize}
  \vspace{-0.1cm}
  \begin{block}{Objetivo de la Sesión}
    Extender la distribución Binomial a ensayos con $k \ge 3$ categorías mutuamente excluyentes, explotar la estructura de covarianzas negativas inducida por la restricción $\sum_{i} X_{i} = n$, y aplicar la distribución condicional multinomial al análisis de tablas de contingencia.
  \end{block}
\end{frame}

% Slide 3: Motivación
\begin{frame}{De la Binomial a la Multinomial: Ensayos con $k$ Categorías}
  \footnotesize
  \begin{motivacion}[¿Por qué generalizar la Binomial?]
    En muchas aplicaciones reales, cada ensayo no tiene solo dos resultados posibles (éxito/fracaso), sino $k \ge 2$ categorías mutuamente excluyentes. Por ejemplo, en una encuesta política los votantes eligen entre 3 candidatos; en genómica, los genotipos pueden ser homocigoto dominante, heterocigoto u homocigoto recesivo; en un dado, hay 6 caras. La distribución \emph{multinomial} generaliza la Binomial a este escenario vectorial.
  \end{motivacion}

  \vspace{0.15cm}
  \pause
  \begin{itemize}\itemsep=0.1cm
    \item \textbf{Vector aleatorio:} $\bm{X} = (X_{1}, X_{2}, \dots, X_{k})$ cuenta cuántos ensayos caen en cada categoría, con $\sum_{i=1}^{k} X_{i} = n$. \pause
    \item \textbf{Parámetros:} Número de ensayos $n$ y probabilidades $p_{1}, p_{2}, \dots, p_{k}$ con $\sum p_{i} = 1$. \pause
    \item \textbf{Marginales binomiales:} Cada $X_{i} \sim \text{Bin}(n, p_{i})$; las dependencias se manifiestan en las covarianzas negativas.
  \end{itemize}
\end{frame}

% Slide 4: Definición Formal
\begin{frame}{Definición Formal de la Distribución Multinomial}
  \footnotesize
  Un vector aleatorio $\bm{X} = (X_{1}, \dots, X_{k})$ sigue una distribución multinomial $\text{Mult}(n, p_{1}, \dots, p_{k})$ con $n$ ensayos y probabilidades $p_{i}$ si su PMF conjunta es: \pause

  \vspace{0.1cm}
  \begin{block}{Función de Masa de Probabilidad Conjunta}
    $$P(X_{1} = x_{1}, \dots, X_{k} = x_{k}) = \dfrac{n!}{x_{1}! \cdots x_{k}!} p_{1}^{x_{1}} \cdots p_{k}^{x_{k}}, \quad \sum_{i=1}^{k} x_{i} = n$$
  \end{block}

  \vspace{0.1cm}
  \pause
  \begin{alertblock}{Interpretación y Parametrización en SciPy}
    \begin{itemize}\itemsep=0.06cm
      \item \textbf{Coeficiente multinomial:} $\dfrac{n!}{x_{1}! \cdots x_{k}!}$ cuenta las configuraciones etiquetadas (cada ensayo es distinguible).
      \item \textbf{Normalización:} $\sum_{\bm{x}} P(\bm{X} = \bm{x}) = (p_{1} + \cdots + p_{k})^{n} = 1$ por el Multinomio de Newton.
      \item \textbf{En SciPy:} \texttt{multinomial.pmf(x, n, p)} con $\bm{x}$ y $\bm{p}$ vectores; \texttt{np.random.multinomial(n, p, size)} para muestreo.
    \end{itemize}
  \end{alertblock}
\end{frame}

% Slide 5: Marginales Binomiales
\begin{frame}{Reducción Marginal: Cada $X_{i} \sim \text{Bin}(n, p_{i})$}
  \footnotesize
  La distribución multinomial es univariadamente binomial: cada componente marginal es una Binomial. La demostración procede sumando sobre las demás componentes. \pause

  \vspace{0.1cm}
  \begin{block}{Teorema de Marginalización}
    $$P(X_{i} = x_{i}) = \binom{n}{x_{i}} p_{i}^{x_{i}} (1 - p_{i})^{n - x_{i}}, \quad X_{i} \sim \text{Bin}(n, p_{i})$$
    La demostración usa el Teorema del Binomio al sumar sobre todas las configuraciones de las demás $k - 1$ componentes.
  \end{block}

  \vspace{0.1cm}
  \pause
  \begin{alertblock}{Esperanza y Varianza Marginales}
    \begin{itemize}\itemsep=0.05cm
      \item $\E[X_{i}] = n p_{i}$
      \item $\Var(X_{i}) = n p_{i} (1 - p_{i})$
      \item Las marginales \emph{no capturan} las dependencias; se requiere la matriz de covarianzas para el análisis conjunto.
    \end{itemize}
  \end{alertblock}
\end{frame}

% Slide 6: Covarianzas Negativas
\begin{frame}{Estructura de Covarianzas: $\cov(X_{i}, X_{j}) = -n p_{i} p_{j}$}
  \footnotesize
  A diferencia de las marginales, las covarianzas conjuntas codifican las dependencias inducidas por la restricción $\sum_{i} X_{i} = n$. \pause

  \vspace{0.1cm}
  \begin{block}{Teorema de Covarianzas Multinomiales}
    $$\cov(X_{i}, X_{j}) = -n p_{i} p_{j}, \quad i \neq j$$
    Representamos $X_{i} = \sum_{r=1}^{n} I_{r}^{(i)}$ (indicadores por ensayo). Por independencia entre ensayos, solo los términos $r = s$ contribuyen, donde $\cov(I_{r}^{(i)}, I_{r}^{(j)}) = -p_{i} p_{j}$ por exclusión mutua.
  \end{block}

  \vspace{0.1cm}
  \pause
  \begin{alertblock}{Matriz de Covarianzas Completa}
    La matriz $\bm{\Sigma}$ tiene diagonal $np_{i}(1-p_{i})$ y fuera de la diagonal $-np_{i}p_{j}$. Es \emph{singular} de rango $k - 1$ porque $\sum_{i} X_{i} = n$ impone una restricción lineal.
  \end{alertblock}
\end{frame}

% Slide 7: Distribución Condicional Multinomial
\begin{frame}{Distribución Condicional: Sub-vector dado Marginal Fijo}
  \footnotesize
  Si $\bm{X} = (X_{1}, \dots, X_{k}) \sim \text{Mult}(n, p_{1}, \dots, p_{k})$ y fijamos $X_{k} = m_{k}$, la distribución del sub-vector $(X_{1}, \dots, X_{k-1})$ es multinomial con $n - m_{k}$ ensayos y probabilidades renormalizadas: \pause

  \vspace{0.1cm}
  \begin{block}{Teorema de Condicional Multinomial}
    $$(X_{1}, \dots, X_{k-1}) \mid (X_{k} = m_{k}) \sim \text{Mult}(n - m_{k}, p_{1}', \dots, p_{k-1}'), \quad p_{i}' = \dfrac{p_{i}}{1 - p_{k}}$$
  \end{block}

  \vspace{0.1cm}
  \pause
  \begin{alertblock}{Caso Particular: Condicional Binomial}
    Para $k = 3$, condicionar en $X_{2} + X_{3} = m$ produce:
    $$X_{2} \mid (X_{2} + X_{3} = m) \sim \text{Bin}(m, p_{2}/(p_{2} + p_{3}))$$
    Esta es la base de la prueba exacta de Fisher en tablas de contingencia.
  \end{alertblock}
\end{frame}

% Slide 8: Aplicaciones
\begin{frame}{Aplicaciones en Ciencia de Datos y Análisis Multivariado}
  \footnotesize
  La distribución multinomial es la columna vertebral del análisis de datos categóricos: \pause

  \vspace{0.15cm}
  \begin{itemize}\itemsep=0.15cm
    \item \textbf{Encuestas de opinión y ciencia política:} Distribución de preferencias entre múltiples candidatos o partidos en una muestra de votantes. \pause
    \item \textbf{Gestión de calidad industrial:} Clasificación de artículos manufacturados en múltiples categorías de calidad (Premium, Estándar, Defectuoso, etc.). \pause
    \item \textbf{Genómica de poblaciones:} Distribución de genotipos bajo equilibrio de Hardy-Weinberg con tres genotipos (homocigoto dominante, heterocigoto, homocigoto recesivo). \pause
    \item \textbf{Procesamiento de lenguaje natural:} Distribución de frecuencias de palabras o categorías gramaticales en un corpus de documentos.
  \end{itemize}
\end{frame}

% Slide 12A-1: Lab Python (Bloque 1)
\begin{frame}[fragile]{Laboratorio en Python: PMF y Marginalización}
  \scriptsize
  Verificación computacional en SciPy (\texttt{03.08\_multinomial\_distribution.py}) de la PMF multinomial, normalización y reducción marginal a Binomial:
  \vspace{0.05cm}
  {\lstset{basicstyle=\fontsize{5pt}{6pt}\selectfont\ttfamily}
  \lstinputlisting[firstline=15, lastline=39]{../../code/03_variables_aleatorias_discretas/03.08_multinomial_distribution.py}}
\end{frame}

% Slide 12A-2: Lab Python (Bloque 2)
\begin{frame}[fragile]{Laboratorio en Python: Covarianzas Negativas}
  \scriptsize
  Simulación Monte Carlo con $250{,}000$ ensayos verificando $\cov(X_{i}, X_{j}) = -n p_{i} p_{j}$:
  \vspace{0.02cm}
  {\lstset{basicstyle=\fontsize{4.5pt}{5.5pt}\selectfont\ttfamily}
  \lstinputlisting[firstline=43, lastline=72]{../../code/03_variables_aleatorias_discretas/03.08_multinomial_distribution.py}}
\end{frame}

% Slide 12B: Lab Python (Bloque 3)
\begin{frame}[fragile]{Laboratorio en Python: Distribución Condicional Multinomial}
  \scriptsize
  Verificación numérica de que $X_{3} \mid (X_{3} + X_{4} = m) \sim \text{Bin}(m, p_{3}/(p_{3} + p_{4}))$:
  \vspace{0.02cm}
  {\lstset{basicstyle=\fontsize{4.5pt}{5.5pt}\selectfont\ttfamily}
  \lstinputlisting[firstline=83, lastline=108]{../../code/03_variables_aleatorias_discretas/03.08_multinomial_distribution.py}}
\end{frame}

% Slide 12C: Salida Terminal
\begin{frame}[fragile]{Laboratorio en Python: Salida en Terminal}
  \scriptsize
  Salida estándar generada al ejecutar el script del laboratorio en Python:
  \vspace{0.02cm}
  \begin{block}{Salida Estándar en Consola}
    \fontsize{4pt}{5pt}\selectfont
    \begin{verbatim}
=== Block 1: Multinomial PMF & Marginalization ===
n=100, p=[0.45, 0.35, 0.20], target=(50, 30, 20)
P(X=(50,30,20)) = 4.79e-03
PMF normalization: 1.000000
Marginal X_1 ~ Bin(100, 0.45): P(X_1=50) theo=0.048152 | joint=0.048152

=== Block 2: Covariance Structure & Negative Correlations ===
Theoretical: E[X]=(50,30,20) | Var=(25,21,16)
  Cov(X_1,X_2)=-15 | Cov(X_1,X_3)=-10 | Cov(X_2,X_3)=-6
Empirical (N=250,000): E[X]=(50.0,30.0,20.0)
  Cov(X_1,X_2)=-14.96 | Cov(X_1,X_3)=-10.12 | Cov(X_2,X_3)=-5.93

=== Block 3: Conditional Distributions ===
Mult(100, [0.4, 0.3, 0.2, 0.1]) | p_cond = 0.6667
Conditional X_3 | (X_3+X_4=30) ~ Bin(30, 0.6667)

Theoretical Bin(30, 0.6667):
  k=18: 0.1101 | k=19: 0.1391 | k=20: 0.1530 | k=21: 0.1457

Empirical conditional (N=500,000, |filter|=43,324):
  k=18: emp=0.1109 | theo=0.1101 | k=20: emp=0.1497 | theo=0.1530

Problem 3.8.9 verification: P(X_3=18 | X_3+X_4=30) = 0.1101
    \end{verbatim}
  \end{block}
\end{frame}

% Slide 13A: Ejercicio Nivel 1 (Enunciado)

% Slide 13B: Ejercicio Nivel 1 (Resolución)

% Slide 14A: Ejercicio Nivel 2 (Enunciado)

% Slide 14B: Ejercicio Nivel 2 (Resolución)

% Slide 15A: Ejercicio Nivel 3 (Enunciado)

% Slide 17: Cierre
\begin{frame}{Cierre de Sección y Síntesis sobre Distribuciones Multivariadas}
  \begin{reflexion}[El Poder de las Categorías Múltiples]
    La distribución multinomial nos enseña que, al extender la Binomial a múltiples categorías, las marginales se preservan (cada $X_{i}$ sigue siendo binomial) pero emergen nuevas dependencias estructuradas: las covarianzas negativas $\cov(X_{i}, X_{j}) = -n p_{i} p_{j}$ codifican la restricción fundamental $\sum X_{i} = n$.
  \end{reflexion}

  \vspace{0.2cm}
  \pause
  \begin{block}{Lecciones Clave de la Modelación Multinomial}
    \begin{itemize}\itemsep=0.08cm
      \item \textbf{Marginales preservadas:} Cada $X_{i} \sim \text{Bin}(n, p_{i})$ individualmente, pero las dependencias son esenciales para el análisis conjunto.
      \item \textbf{Covarianzas negativas:} Reflejan la restricción $\sum X_{i} = n$; un exceso en $i$ implica déficit en alguna $j$.
      \item \textbf{Condicional multinomial:} Bajo marginales fijos, las distribuciones colapsan a multinomiales con menos categorías, base de la prueba exacta de Fisher.
    \end{itemize}
  \end{block}
\end{frame}

% Slide 18: Perspectiva Modular
\begin{frame}{Perspectiva Modular --- El Próximo Reto}
  \scriptsize
  \begin{retoclase}[Para la próxima sesión: Distribución Normal y Aproximación Continua]
    Hasta ahora hemos estudiado distribuciones de probabilidad para variables aleatorias \emph{discretas} que toman valores enteros. Pero muchos fenómenos naturales (altura, peso, errores de medición) varían continuamente y siguen una distribución en forma de campana. Introduciremos la \emph{Distribución Normal} y su papel como límite continuo de la Binomial (De Moivre-Laplace).
  \end{retoclase}

  \vspace{0.08cm}
  \pause
  \begin{alertblock}{Preparación Recomendada}
    \begin{itemize}\itemsep=0.06cm
      \item Repasa el concepto de función de densidad de probabilidad (PDF) y su diferencia con la PMF.
      \item Reflexiona sobre la conexión asintótica entre la Binomial y la Normal cuando $n$ crece manteniendo $p$ fija.
    \end{itemize}
  \end{alertblock}
\end{frame}

\end{document}
